% ---------------------------------------------------------------------------
% Preambule -- PFE Toulouse INP-ENSEEIHT / MF2E
%
% Contraintes de forme imposees (student_context/Consignes pour les PFE2526.pdf) :
%   - coeur <= 30 pages, POLICE 12, figures comprises, hors annexes
%   - annexes <= 20 pages
%   - resume <= 2 pages en francais ET en anglais
%   - page de garde, remerciements, sommaire pagine, table des acronymes, biblio en fin
% Les pages de garde / biblio / sommaire / resume ne comptent pas dans les 30.
% ---------------------------------------------------------------------------

% --- Langue et encodage ----------------------------------------------------
\usepackage[T1]{fontenc}
\usepackage[utf8]{inputenc}
% `english` doit etre charge aussi : le resume anglais l'appelle via
% \begin{otherlanguage}{english}. Le DERNIER de la liste est la langue principale.
\usepackage[english,french]{babel}
\frenchsetup{SmallCapsFigTabCaptions=false}
\usepackage{csquotes}

% --- Fontes ----------------------------------------------------------------
% newtx (Times) : lisible en 12 pt et nettement plus compact que Latin Modern,
% ce qui compte quand 30 pages est une contrainte dure et non un objectif.
\usepackage{newtxtext}
\usepackage{newtxmath}
\usepackage{microtype}

% --- Mise en page ----------------------------------------------------------
\usepackage[a4paper,top=2.5cm,bottom=2.5cm,left=2.5cm,right=2.5cm]{geometry}
\usepackage{setspace}
\onehalfspacing
\usepackage{fancyhdr}
\pagestyle{fancy}
\fancyhf{}
\fancyhead[L]{\footnotesize\nouppercase{\leftmark}}
\fancyhead[R]{\footnotesize\thepage}
\renewcommand{\headrulewidth}{0.4pt}
\fancypagestyle{plain}{\fancyhf{}\fancyfoot[C]{\footnotesize\thepage}%
  \renewcommand{\headrulewidth}{0pt}}

% --- Tableaux et figures ---------------------------------------------------
\usepackage{graphicx}
\usepackage{booktabs}
\usepackage{tabularx}
\usepackage{multirow}
\usepackage{longtable}
\usepackage{array}
\usepackage[font=small,labelfont=bf,width=0.92\textwidth]{caption}
\usepackage{float}

% --- Mathematiques ---------------------------------------------------------
\usepackage{amsmath}
% newtxmath definit deja \Bbbk ; amssymb le redefinirait et LaTeX s'arrete.
\let\Bbbk\relax
\usepackage{amssymb}
\usepackage{mathtools}

% --- Nombres et unites -----------------------------------------------------
% Convention francaise : virgule decimale, espace insecable fine aux milliers.
\usepackage{siunitx}
\sisetup{
  locale = FR,
  group-separator = {\,},
  group-minimum-digits = 4,
  output-decimal-marker = {,},
  detect-all
}
\usepackage{textcomp}
% LE SYMBOLE EURO. Il apparait des centaines de fois dans ce memoire, il doit donc etre juste.
% newtx ne fournit pas d'euro en TS1 : \texteuro y declenche un "Faking \texteuro", et pdflatex
% fabrique alors le glyphe a la volee dans une fonte bitmap Computer Modern. Le resultat est
% visiblement etranger au texte -- graisse, chasse et dessin ne sont pas ceux de la page.
% `eurosym` fournit le glyphe officiel en Type 1, vectoriel, et s'accorde a un romain classique.
% NE PAS ecrire \DeclareSIUnit{\euro}{\text{\euro}} -- la definition serait recursive.
\usepackage{eurosym}
\newcommand{\eur}{\euro}
\DeclareSIUnit{\ha}{ha}
\DeclareSIUnit{\etp}{ETP}
\DeclareSIUnit{\tonne}{t}
% L'unite siunitx s'appelle \euros et NON \EUR : eurosym definit deja un \EUR{<montant>} qui
% prend un argument, et \DeclareSIUnit ne le remplace pas -- le symbole avalait alors
% l'accolade suivante ("Argument of \EUR has an extra }"). Le \text{} remet la macro en mode
% texte, siunitx composant ses unites en mode mathematique.
\DeclareSIUnit{\euros}{\text{\euro}}

% --- Listes ----------------------------------------------------------------
\usepackage{enumitem}
\setlist{noitemsep,topsep=4pt}

% --- Encadres reflexifs ----------------------------------------------------
% La reflexivite est TISSEE dans chaque chapitre, pas reportee a la fin : chaque
% chapitre se clot sur un encadre. Style sobre, filet a gauche, pas de couleur
% de fond -- il doit se lire comme une respiration, pas comme une boite a outils.
\usepackage[most]{tcolorbox}
\newtcolorbox{reflexion}[1][]{%
  enhanced, breakable, sharp corners,
  colback = black!2, colframe = black!55,
  boxrule = 0pt, leftrule = 1.6pt,
  left = 8pt, right = 8pt, top = 6pt, bottom = 6pt,
  fontupper = \small,
  title = {\itshape Retour réflexif},
  coltitle = black, colbacktitle = black!2,
  fonttitle = \small\itshape\bfseries,
  attach title to upper = {\\[2pt]},
  #1
}

% Encadre pour les mises en garde de lecture (les "pieges" de 04-vigilance.md
% qui doivent voyager avec le chiffre qu'ils qualifient).
\newtcolorbox{vigilance}[1][]{%
  enhanced, breakable, sharp corners,
  colback = white, colframe = black!70,
  boxrule = 0.6pt,
  left = 8pt, right = 8pt, top = 6pt, bottom = 6pt,
  fontupper = \small,
  #1
}

% --- Bibliographie ---------------------------------------------------------
\usepackage[backend=biber,style=authoryear,maxcitenames=2,maxbibnames=99,
            giveninits=true,uniquename=init,sorting=nyt]{biblatex}
\addbibresource{biblio.bib}

% --- Liens -----------------------------------------------------------------
\usepackage[hidelinks,breaklinks]{hyperref}
\usepackage{cleveref}
\crefname{chapter}{chapitre}{chapitres}
\crefname{section}{section}{sections}
\crefname{figure}{figure}{figures}
\crefname{table}{tableau}{tableaux}
\crefname{equation}{équation}{équations}

% --- Acronymes -------------------------------------------------------------
\usepackage[acronym,nomain,nonumberlist,nopostdot]{glossaries}
\makenoidxglossaries

% --- Titres de chapitre sobres --------------------------------------------
\usepackage{titlesec}
\titleformat{\chapter}[display]
  {\normalfont\LARGE\bfseries}{\normalsize\mdseries\textsc{\chaptertitlename~\thechapter}}{6pt}{\LARGE}
\titlespacing*{\chapter}{0pt}{-10pt}{24pt}

% --- Raccourcis ------------------------------------------------------------
\newcommand{\mosaica}{\textsc{Mosaica}}
\newcommand{\gams}{\textsc{Gams}}
\newcommand{\highs}{\textsc{HiGHS}}
\newcommand{\pad}{\textsc{pad}}

% Marqueur de section non redigee. Visible dans le PDF (on ne rend pas un
% memoire ou il en reste) et signale dans le .log pour pouvoir les compter.
\newcommand{\arediger}[1]{%
  \typeout{A REDIGER : #1}%
  \par\medskip\noindent\fbox{\parbox{\dimexpr\linewidth-2\fboxsep-2\fboxrule}{%
    \small\itshape À rédiger --- #1}}\par\medskip}

% Compte de pages du coeur : le compteur arabe redemarre a 1 au chapitre 1 et
% \pageref{fin-coeur} donne donc directement le nombre de pages du coeur. Il est
% ecrit dans le .log a chaque compilation pour que la contrainte des 30 pages
% soit verifiee mecaniquement et non a l'oeil.
\newcommand{\annoncecoeur}{%
  \typeout{^^J===== PAGES DU COEUR : \pageref{fin-coeur} / 30 =====^^J}}
